\section*{\S4 Medlemsmöte}
\addcontentsline{toc}{section}{\S4 Medlemsmöte}

\subsection*{\S4.1 Befogenhet}
\phantomsection
\addcontentsline{toc}{subsection}{\S4.1 Befogenhet}

Medlemsmötet är SVETT:s högsta beslutande organ.

\subsection*{\S4.2 Sammansättning}
\addcontentsline{toc}{subsection}{\S4.2 Sammansättning}

Medlemsmötet utgörs av SVETT:s medlemmar.

\subsection*{\S4.3 Sammanträde}
\addcontentsline{toc}{subsection}{\S4.3 Sammanträde}

Medlemsmöte ska hållas minst en gång per termin, dock inte under helgdag, lov eller
tentamensperiod.

\subsection*{\S4.4 Sammankallande}
\addcontentsline{toc}{subsection}{\S4.4 Sammankallande}

Medlemsmöte sammanträder på kallelse av styrelsen. Följande har rätt att skriftligen
begära att styrelsen sammankallar medlemsmöte:
\begin{enumerate}[label=\roman*)]
	\item Datasektionens styrelse,
	\item minst tio (10) medlemmar som står bakom en sådan begäran.
\end{enumerate}

\subsection*{\S4.5 Tidpunkt}
\addcontentsline{toc}{subsection}{\S4.5 Tidpunkt}

Om en begäran enligt \S4.4 har inkommit till styrelsen ska medlemsmöte hållas inom
femton (15) studiedagar.

\subsection*{\S4.6 Kallelse}
\addcontentsline{toc}{subsection}{\S4.6 Kallelse}

Kallelse till medlemsmöte, inklusive förslag till föredragningslista och bilagor, ska
vara medlemmarna tillhanda senast fem (5) dygn före mötet. Kallelsen ska skickas via
e-post till samtliga SVETT:s medlemmar.

\subsection*{\S4.7 Beslutsmässighet}
\addcontentsline{toc}{subsection}{\S4.7 Beslutsmässighet}

Medlemsmötet är beslutsmässigt när:
\begin{enumerate}[label=\roman*)]
	\item kallelse har gått ut enligt \S4.6,
	\item minst tre (3) medlemmar, utöver styrelsen, är närvarande under hela mötet,
	\item Datasektionens styrelse har fått kännedom om mötet.
\end{enumerate}

\subsection*{\S4.8 Röstning}
\addcontentsline{toc}{subsection}{\S4.8 Röstning}

Beslut fattas med enkel majoritet. Vid lika röstetal har SVETT:s ordförande utslagsröst,
om ingen medlem yrkar på bordläggning. Undantag från dessa regler gäller vid
personval enligt \S4.9, beslut om stadgeändring enligt \S8.1 samt beslut om
reglementsändring enligt \S8.2. Varje närvarande medlem har en (1) röst.
Röstning genom fullmakt är inte tillåten.

\subsection*{\S4.9 Personval}
\addcontentsline{toc}{subsection}{\S4.9 Personval}

Personval sker med enkel majoritet. Valet sker genom sluten omröstning om det finns
flera kandidater till en post, eller om minst en (1) närvarande medlem begär det. Vid
lika röstetal bordläggs ärendet till nästa medlemsmöte, vilket ska sammankallas inom
femton (15) studiedagar. Personval sker enligt \S7.1.

\subsection*{\S4.10 Protokoll}
\addcontentsline{toc}{subsection}{\S4.10 Protokoll}

Protokoll ska föras vid medlemsmöte och justeras av två (2) personer som väljs vid
mötet. Protokollet ska göras tillgängligt på internet så snart som möjligt.

\subsection*{\S4.11 Motioner}
\addcontentsline{toc}{subsection}{\S4.11 Motioner}

Varje medlem har rätt att lämna motion till medlemsmöte. Motionen ska vara
styrelsen tillhanda senast tio (10) dygn före mötet och behandlas av SVETT:s styrelse.

\subsection*{\S4.12 Övriga frågor}
\addcontentsline{toc}{subsection}{\S4.12 Övriga frågor}

Under punkten ``Övriga frågor'' kan inga beslut fattas.

\subsection*{\S4.13 Föregående verksamhetsår}
\addcontentsline{toc}{subsection}{\S4.13 Föregående verksamhetsår}

På det första medlemsmötet under höstterminen ska följande ärenden avseende
föregående verksamhetsår behandlas:
\begin{enumerate}[label=\roman*)]
	\item verksamhetsberättelse,
	\item ekonomisk berättelse,
	\item styrelsens ansvarsfrihet.
\end{enumerate}
